\documentclass[11pt]{article}
\usepackage[onepage,ko]{mas1004}

\reporttitle{과제 2: 짧은 보고서}
\studentname{김윤서}
\studentid{20251129}

\begin{document}
\makeheader

\section*{I. 실험 목적 및 배경}

Data \& AI 강의에서 배운 인공지능 모델을 이용해 서로 다른 5개 도시(서울, 부산, 도쿄, 오사카, 타이베이)의 길거리 사진을 보고 도시를 맞추는 과제를 진행했습니다.

가장 먼저 제공된 기본 코드(Baseline)로 1차 학습을 돌려보았는데, 전체 정확도는 49.56\%가 나왔습니다. 도시별로 결과를 확인해 보았더니 사진 수가 많은 도쿄는 62.82\%의 정확도를 기록한 반면, 사진 수가 적고 비교적 어두운 사진이 많은 오사카(34.02\%)와 부산(38.56\%)의 정확도가 특히 낮음을 알 수 있었습니다. 이번 과제에서는 한 번의 시도에 한 번의 조건만 변경하는 것을 원칙으로 데이터 증강(Data Augmentation)을 활용해 총 3차에 걸쳐 실험해 보았습니다.

\section*{II. 세 번의 실험과 결과 비교}

\begin{enumerate}[label=\arabic* 서비스 학습:, leftmargin=1.5em]
  \item[\bfseries ① 1차 학습 :] 주어진 코드 그대로 실행하여 도쿄나 서울과 같은 깔끔한 사진에 익숙해지고 비교적 그렇지 않은 도시는 정확도가 좋지 않았습니다.
  \item[\bfseries ② 2차 학습 :] 가장 낮은 정확도를 기록한 오사카의 정확도를 올리는 것을 목표로, 기존에 설정되어 있던 자르기 범위를 줄이고(Crop 0.5), 사진에 흐림 효과(GaussianBlur)와 강한 색상 변형을 하여 다시 학습했습니다. 결과 오사카의 정확도는 6.25\%p나 상승한 40.27\%를 기록했습니다. 하지만 가장 높은 정확도를 기록했던 도쿄의 정확도가 떨어졌습니다. 전체 정확도는 49.56\%로 1차 학습과 동일한 결과를 확인했습니다.
  \item[\bfseries ③ 3차 학습 :] 2차 학습 때 추가한 블러 효과가 1차 학습에 효과적이었던 특징들을 가려 정확도가 오히려 떨어졌다고 예상하고 블러 효과를 뺀 뒤 색상 변화를 일부 조정했습니다. 1차, 2차 학습과 비교하여 도쿄(61.99\%)와 서울(57.78\%)의 정확도는 비슷한 수준으로 회복했지만 오사카의 정확도는 다시 떨어졌습니다.
\end{enumerate}

\section*{III. 과제를 통해 느낀 점과 배운 점}

\begin{itemize}
  \item 각 도시마다 데이터의 특징이 다르기 때문에 모든 도시 구분의 정확도를 올리는 증강 기법을 찾기 어려웠습니다. 2차 학습에서 블러 처리와 색상 변형으로 흐릿한 데이터의 오사카 구분 정확도를 높인 반면 도쿄와 서울의 정확도가 떨어졌고 3차 학습에서 이 처리를 줄이니 다시 반대로 돌아갔습니다. 이를 통해 도시별로 데이터에 따라 필요한 데이터 변형 기법이 다른 트레이드오프(Trade-off) 현상을 직접 확인해 볼 수 있었습니다. 추후 정확도 상승을 위해서는 각 데이터별 특성을 파악하고 그에 맞는 변형 기법을 적용할 필요가 있다고 생각됩니다.
  \item 이미지만을 변형해 학습한 결과로는 정확도가 3차의 학습 모두에서 절반을 넘지 못하고 49.56\%에 머물렀습니다. 특정 도시의 정확도가 높아져도 전체 도시 모두를 고르게 높이기 어려우니 모델 파라미터를 통한 수정에 한계가 있음을 알게 되었습니다. 이를 해결하기 위해 직접 데이터 개수를 고르게 맞추거나 분류에 오류가 생길 수 있는 데이터들을 고르는 작업이 필요한 것으로 예상됩니다. 하지만 4차, 5차 추가 학습까지 진행해 보며 전체 정확도가 어떻게 같은 수치로 유지되는지 더 학습과 결과 도출을 통해 확인해 보고 싶기도 합니다.
\end{itemize}

\figurepages

\begin{figure}[htbp]
  \centering
  \includegraphics[width=0.85\linewidth]{commons_7266637.jpg}
  \includegraphics[width=0.85\linewidth]{commons_9476105.jpg}
  \caption{선명한 건물 구조와 도로 표지판이 특징인 서울 및 도쿄 길거리 데이터 예시}
\end{figure}

\clearpage

\begin{figure}[htbp]
  \centering
  \includegraphics[width=0.85\linewidth]{commons_2691073.jpg}
  \includegraphics[width=0.55\linewidth]{commons_4000589.jpg}
  \caption{조명 변형 및 노이즈가 비교적 크게 관찰되는 오사카 및 부산 길거리 데이터 예시}
\end{figure}

\end{document}
